% =============================================================================
% TRACTABILITY AND LOCAL APPROXIMATION (3-4 pages)
%
% Spec requirements:
%   - Intractability result (two arguments with different epistemic status)
%   - Definition 7: Local Volume Estimator
%   - Proposition 3: Self-Balancing Survives Approximation
%   - Observability channels (individual + coalition signals)
%   - Remark on estimator failure modes (4 cases)
%   - The feedback loop (re-derive from vol(G) objective)
%   - Computational cost analysis
%
% Reference material:
%   - Predraft Equation 5 (loop structure, needs re-derivation)
%   - Predraft Equations 22-25 (trust model concepts, re-derive notation)
%   - Blog: limitsofthought.md (computational bounds, motivational)
%   - Dyer & Frieze (1988), Klyubin et al. (2005)
% =============================================================================

\section{Tractability and Local Approximation}
\label{sec:tractability}

The definitions and propositions of the preceding sections assume exact access
to $\vol(G)$. This section confronts the fact that no implementable system will
have such access. The argument proceeds in three stages: first, we establish
that exact computation of $\vol(G)$ is intractable; second, we identify the
margin condition under which a local finite-difference estimator preserves
sign-correctness; third, we derive the optimization loop that a GFM actor would
execute in practice.

% ---------------------------------------------------------------------------
% The Intractability Result
% Two separate arguments: formal (#P-hard, external) and motivating
% (combinatorial interaction effects, original, informal)
% ---------------------------------------------------------------------------

\subsection{Intractability of Exact GFM}

Two arguments establish that exact GFM is computationally infeasible. They have
different epistemic status and we present them separately.

\paragraph{Formal result (external).} \citet{dyer1988complexity} showed that
computing the volume of a convex body in $d$ dimensions is \#P-hard. If the
joint goal space $G$ can be represented as a polytope in capability space (or
any region whose volume computation is at least as hard), then exact $\vol(G)$
is intractable. This is a known result being applied to the GFM setting, not a
new proof. The result holds even for convex bodies; the joint goal space $G$ is
generically non-convex (the union in Equation~\eqref{eq:joint_goal_space} does
not preserve convexity), which makes the problem at least as hard.

\paragraph{Motivating observation (original, informal).} The joint goal space
$G = \bigl(\bigcup_k \Gind{k}\bigr) \cup \Gcoop$ includes cooperative
capabilities that arise from multi-agent interaction. Each subset $S \subseteq
\{a_1, \ldots, a_n\}$ with $|S| \geq 2$ can contribute cooperative capabilities
to $\Gcoop$ that no proper subset of $S$ can realize. The number of such
subsets is $2^n - n - 1$. Not every subset generates unique capabilities, but
the combinatorial growth is suggestive: for a population of 60 agents, there
are more than $10^{18}$ possible coalitions (already exceeding the number of
seconds since the Big Bang), and the space of interaction effects among those
coalitions dwarfs the individual capability sets. This does not constitute a
formal reduction from a known hard problem to the representation of $G$, but it
motivates the expectation that representing and measuring the cooperative
component of $G$ grows at least combinatorially with population size. A formal
model of why joint goal-space complexity scales this way remains an open
problem.

\paragraph{Combined conclusion.} Even under optimistic assumptions about the
structure of $\G$ (convexity, bounded dimension, well-behaved measure), exact
$\vol(G)$ computation is at minimum \#P-hard. The combinatorial growth of
interaction effects makes this worse in practice. A GFM actor cannot compute
$\vol(G)$; it must work with an approximation.

% ---------------------------------------------------------------------------
% Definition 7: Local Volume Estimator
% ---------------------------------------------------------------------------

\begin{definition}[Local Volume Estimator]
\label{def:local_volume_estimator}
For a proposed action $\pi$, define the \emph{local volume difference}:
\begin{equation}
\label{eq:local_volume_diff}
\Deltahat(\pi) = \Vhat(G \mid \mathrm{do}(\pi)) - \Vhat(G \mid \mathrm{skip}(\pi))
\end{equation}
where $\Vhat$ is an estimate of $\vol(G)$. This requires only that $\vol(G)$ is
    \emph{measurable}, not that it is differentiable or that $\G$ has smooth
    structure.
\end{definition}

The key insight is that a GFM actor does not need to compute $\vol(G)$. It
needs to determine $\sign(\Delta \vol(G) \mid \pi)$: whether a proposed action
expands or contracts the joint capability space. This is a strictly weaker
requirement. Computing volume requires solving a \#P-hard problem; determining
the sign of a volume change requires only that the estimator's error is bounded
below the true signal magnitude.

Definition~\ref{def:local_volume_estimator} does \emph{not} require a metric on
$\G$, a notion of direction in capability space, or continuity of $\vol(G)$
with respect to actions. It requires only that capabilities can be added to or
removed from $G$ and that these changes are observable through the channels
described below. The estimator operates on discrete changes (capabilities
gained or lost), not on gradients through a continuous space.

\paragraph{Observability channels.} The estimator $\Vhat$ draws on two distinct
signal types with different coverage and reliability.

\emph{Individual capability signals.} For each observed agent $a_k$, the GFM
actor can solicit or observe accept/reject responses $R_k$ indicating whether
the agent's individually observable capability set expanded or contracted.
These responses are weighted by the trust factor $T_k$ from the observable goal
model (Definition~\ref{def:observable_goal_model}) and aggregated across the
observation set. This channel covers changes to $\bigcup_k \Gind{k}$ well: when
an action directly affects an agent's capabilities, that agent can report the
change, and the estimator registers it with weight proportional to trust.

\emph{Coalition capability signals.} Changes to $\Gcoop$ are not directly
observable from individual agents' reports. A cooperative capability can appear
or disappear even when no single agent reports a unilateral change; a
coordination mechanism can dissolve without any individual losing personal
capabilities. The estimator can partially observe $\Gcoop$ through indirect
signals: new coordination patterns between agents (observed behaviorally),
agents reporting capabilities conditional on cooperation with specific
partners, and changes in the structure of agent interactions. But this channel
is noisier, slower, and less reliable than individual signals.

\paragraph{Critical assumption.} $\Vhat$ is a better estimator of
$\vol\bigl(\bigcup_k \Gind{k}\bigr)$ than of $\vol(G)$. The cooperative
component $\Gcoop$ is partially observable but systematically underweighted.
This is the primary implementation gap in the framework: the local estimator
provides a lower bound on the information available about $\vol(G)$, and its
coverage of cooperative capabilities is where implementations must invest most
heavily to close the gap between the estimator's reach and the objective's
scope.

% ---------------------------------------------------------------------------
% Proposition 3: Self-Balancing Survives Approximation
% ---------------------------------------------------------------------------

\begin{proposition}[Sign-Correctness Decomposition]
\label{prop:sign_correctness}
Decompose $\Delta \vol(G)$ into the individual-capability component $\Delta
    \vol\bigl(\bigcup_k \Gind{k}\bigr)$ and the cooperative component $\Delta
    \vol(\Gcoop)$. The local volume estimator $\Deltahat(\pi)$ observes the
    first component (through agent reports) and partially observes the second
    (through coalition signals). The estimator's sign matches the true sign,
    $\sign(\Deltahat(\pi)) = \sign(\Delta \vol(G) \mid \pi)$, when all three of
    the following conditions hold:
\begin{enumerate}
\item The actor observes a representative sample of affected agents.
\item The actor's trust model is accurate enough to distinguish genuine feedback from adversarial noise.
\item \textbf{Individual-level dominance:} There exists $\epsilon < 1$ such that
\begin{equation}
\label{eq:sign_preservation}
\bigl|\Delta \vol(\Gcoop) - \widehat{\Delta \vol}(\Gcoop)\bigr| < (1 - \epsilon) \cdot \bigl|\Delta \vol\bigl(\textstyle\bigcup_k \Gind{k}\bigr)\bigr|
\end{equation}
\end{enumerate}
\end{proposition}

The proposition is a \emph{decomposition}, not a derivation of tractability.
Condition~\eqref{eq:sign_preservation} is a margin condition: it states that
the estimation error on the cooperative component is smaller than the observed
individual-level signal. Given that margin, sign-correctness follows
immediately---the contribution is in identifying \emph{which actions satisfy
the margin condition naturally} and which do not, thereby telling implementers
where the estimator is reliable and where it requires additional observation
channels.

\emph{Verification.} The self-balancing property
(Proposition~\ref{prop:self_balancing}) is structural: it follows from measure
monotonicity regardless of whether the estimator can observe the full change.
The decomposition bridges from this structural property to the estimator's
output. The estimator observes the individual-capability component (up to
sampling error controlled by conditions 1--2) and partially observes the
cooperative component. Under condition~3, the estimation error on the
cooperative component cannot flip the sign of the observed individual-level
signal. Since the individual signal preserves the sign of the true
individual-level change (by conditions 1--2), the estimator's sign matches the
sign of the true $\Delta \vol(G)$.

\paragraph{When condition 3 holds naturally.} The actions an alignment
framework most needs to evaluate---direct harm, resource destruction, coercion,
capability expansion through education or tool-building---produce large
individual-capability signals across many agents. An action that destroys an
agent's resources visibly contracts $\Gind{k}$; an action that teaches a new
skill visibly expands it. For these actions, the individual-level change
\emph{is} the dominant effect, and condition~3 holds with large margin. The
estimator is strongest exactly where alignment matters most: actions whose
primary mechanism is changing what individual agents can do.

\paragraph{When condition 3 fails.} Actions whose impact is primarily on
coalition structure (dissolving a coordination mechanism, introducing distrust
between allies, or enabling new multi-agent capabilities) while leaving
individual capability sets largely unchanged. These actions change
$\vol(\Gcoop)$ without producing proportionate individual-level signals. They
are real and important, but they are a narrower class of actions than the
individual-capability cases. The paper identifies them as the estimator's known
boundary rather than treating them as representative of all estimation
failures.

% ---------------------------------------------------------------------------
% Remark: Estimator failure modes
% ---------------------------------------------------------------------------

\begin{remark}[Estimator Failure Modes]
\label{rem:failure_modes}
Proposition~\ref{prop:sign_correctness} covers the most critical actions for
    alignment. When its assumptions do not hold, the following characterization
    serves as a threat model telling implementers where to invest in better
    observation, not where to give up.

\begin{enumerate}
\item \textbf{Full individual observation, no cooperative change.} The estimator is sign-correct: it sees the complete relevant change. This covers the common case of actions that directly expand or contract agents' personal capabilities.

\item \textbf{Partial individual observation, no cooperative change.} Sign-correctness depends on the sampling model, which is an implementation design choice. Unbiased sampling of affected agents improves with coverage; adversarially selected observation sets can mislead. Condition~1 of Proposition~\ref{prop:sign_correctness} controls this case.

\item \textbf{Mixed individual and cooperative changes.} The estimator is unreliable when the cooperative-level change dominates the individual-level signal. This is a structural limitation; improving the coalition observation channel is the fix, not better sampling of individuals.

\item \textbf{Pure coalition-level changes.} The estimator is blind. An action that restructures cooperation patterns without affecting any individual's capability set produces no signal in the individual channel. This is the known boundary.
\end{enumerate}

What the Remark gives an implementer: a graduated threat model. The self-balancing property of Proposition~\ref{prop:self_balancing} holds over all of $\vol(G)$, the Remark characterizes how much of that property a given implementation can enforce, and where to extend coverage. Cases~1--2 are addressable through sampling design. Case~3 requires investment in coalition-level observation. Case~4 identifies the structural frontier where new measurement channels are needed.
\end{remark}

% ---------------------------------------------------------------------------
% The Feedback Loop
% Re-derive from vol(G) objective using Definitions 1-6
% ---------------------------------------------------------------------------

\subsection{The Optimization Loop}
\label{subsec:optimization_loop}

Given the local volume estimator, we can now derive the optimization loop a GFM
actor executes at each decision point. The loop re-derives the structure of the
GFM objective (Equation~\eqref{eq:gfm_objective}) in implementable terms, using
the definitions established in
Sections~\ref{sec:definitions}--\ref{sec:self_balancing}.

At each step, the actor:

\begin{enumerate}
\item \textbf{Propose.} Select a candidate action $\pi$ from the action space, informed by the current observable goal models $\{M_k(G)\}$ (Definition~\ref{def:observable_goal_model}). The proposal targets agents whose capability sets the actor believes can be expanded.

\item \textbf{Estimate.} Compute the local volume difference $\Deltahat(\pi)$ (Definition~\ref{def:local_volume_estimator}) by aggregating accept/reject signals $R_k$ from observed agents, each weighted by trust factor $T_k$:
\begin{equation}
\label{eq:aggregate_signal}
\Deltahat(\pi) \;=\; \sum_{k \in \mathcal{O}} T_k \cdot R_k(\pi)
\end{equation}
where $\mathcal{O}$ is the actor's observation set, $R_k(\pi) \in \{-1, 0, +1\}$ represents agent $a_k$'s reported capability change under $\pi$, and $T_k \in [0, 1]$ is the trust factor from Definition~\ref{def:observable_goal_model}. The sign of $\Deltahat(\pi)$ determines whether $\pi$ is classified as expanding or contracting.

\item \textbf{Act or abstain.} Execute $\pi$ if $\Deltahat(\pi) > 0$ (estimated expansion). Abstain or seek alternatives if $\Deltahat(\pi) \leq 0$ (estimated contraction or neutral). In the marginal case where $|\Deltahat(\pi)|$ is small, the actor may defer pending additional evidence.

\item \textbf{Update models.} After observing the consequences of $\pi$ (or its absence), update each observable goal model $M_k(G)$ using the observed changes in agent $a_k$'s capability set:
\begin{equation}
\label{eq:model_update}
M_k'(G) \;=\; M_k(G) + T_k \cdot (1 - T_s) \cdot f(R_k \mid M_k(G)) \cdot \mathbf{1}\bigl[\Delta E(R_k, M_k(G)) > \theta\bigr]
\end{equation}
where $(1 - T_s)$ is the \emph{learning rate} derived from the actor's self-trust $T_s \in [0,1]$: when self-trust is low (bootstrapping phase), the learning rate is high and the model admits many updates; when self-trust is high (stable phase), the learning rate is low and the model resists small perturbations. The term $f(R_k \mid M_k(G))$ is the update direction derived from comparing agent $a_k$'s response $R_k$ against the model's prediction.
\end{enumerate}

\paragraph{Activation energy.} The indicator $\mathbf{1}[\Delta E > \theta]$ in
Equation~\eqref{eq:model_update} implements an \emph{activation energy
threshold}: the world model updates only when the evidence from a response
exceeds a threshold $\theta$. This is a key design parameter for adversarial
robustness, not a minor implementation detail. Without it, an adversary could
incrementally shift the actor's world model through a sequence of small,
individually sub-threshold manipulations: each interaction nudges the model
slightly, and over many interactions the cumulative drift is large. The
activation energy threshold $\theta$ prevents this incremental gaslighting by
requiring that any single piece of evidence must clear a minimum bar before it
can alter the actor's beliefs. Small perturbations are absorbed; only signals
with sufficient magnitude propagate into the world model.

The threshold $\theta$ creates a tradeoff: too high, and the actor becomes
unresponsive to genuine evidence; too low, and it becomes vulnerable to
manipulation. The correct value of $\theta$ is context-dependent and
constitutes one of the primary tuning parameters for a GFM implementation. In
environments with high adversarial pressure, $\theta$ should be large; in
cooperative environments with trusted agents, it can be smaller.

\paragraph{Trust model interface.} The trust factor $T_k$ and self-trust $T_s$
in Equations~\eqref{eq:aggregate_signal}--\eqref{eq:model_update} are learned
quantities, updated through interaction. New agents enter with a low initial
trust $T_k$ during a \emph{cooling period} $\tau$, during which the actor
accumulates evidence about the agent's reliability before weighting its reports
heavily. $T_k$ increases as agent $a_k$'s reported capability changes are
confirmed by independent observation, and decreases when reports diverge from
observed reality. The full derivation of the trust update dynamics is deferred
to Appendix~\ref{app:trust_model}; here we note only that the trust model is
the mechanism by which condition~2 of Proposition~\ref{prop:sign_correctness}
is satisfied in practice.

% ---------------------------------------------------------------------------
% Computational Cost
% ---------------------------------------------------------------------------

\subsection{Computational Cost}
\label{subsec:computational_cost}

The local volume estimator reduces the computational requirements of GFM from
intractable (exact $\vol(G)$ over all agents) to practical (sign estimation
over observed agents).

\paragraph{Scaling.} The per-decision cost of the optimization loop scales with
$|\mathcal{O}|$, the number of observed agents, not with the total population
$n$. Each decision requires: proposing an action (dependent on the actor's
policy complexity), soliciting or observing responses from agents in
$\mathcal{O}$ (linear in $|\mathcal{O}|$), aggregating the trust-weighted
signals (linear in $|\mathcal{O}|$), and updating the goal models $M_k(G)$ for
affected agents (linear in the number of agents whose models change). The
dominant cost is maintaining and updating the observable goal models, which is
bounded by $|\mathcal{O}|$ per step.

\paragraph{Comparison.} This computational profile is comparable to a
contextual multi-armed bandit with structured feedback. The actor selects an
action (arm), observes a structured response from the environment
(accept/reject signals from multiple agents rather than a scalar reward), and
updates its model of the environment. The key structural difference from a
standard bandit is that the feedback is multi-dimensional and trust-weighted:
each agent provides an independent signal, and the aggregation incorporates the
actor's learned assessment of each signal's reliability. The resulting problem
is harder than a scalar bandit but remains tractable---the per-step cost is
polynomial in $|\mathcal{O}|$ and does not depend on the size of the full
capability space $\G$.

\paragraph{What is lost.} The local estimator cannot detect changes to $\Gcoop$
that leave individual capability sets unchanged
(Remark~\ref{rem:failure_modes}, case~4). It systematically underweights
cooperative capabilities. And it provides no guarantee on the \emph{magnitude}
of $\Delta \vol(G)$---only on its sign. A GFM actor operating through the local
estimator may correctly classify actions as expanding or contracting without
knowing how much expansion or contraction occurs. This limits the actor's
ability to compare two expanding actions and select the better one, reducing
the optimization from ``maximize $\vol(G)$'' to ``take expanding actions and
avoid contracting ones.'' The magnitude estimation problem is an open direction
for future work.

\paragraph{What improves with capability.} The limitations above are bounded by
the actor's observational reach and world-model quality, both of which improve
with capability. The three conditions of
Proposition~\ref{prop:sign_correctness} are all monotonically strengthened by
increased resources. Condition~1 (representative sample) improves with
observational reach: a more capable actor maintains a larger, better-calibrated
$\mathcal{O}$, covering more of the population affected by each action.
Condition~2 (trust model accuracy) improves with compute: the gradient-tracking
dynamics of Appendix~\ref{app:trust_model} converge faster and track behavioral
change more reliably with more data and finer-grained agent models. Condition~3
(sign-preservation) improves with world-model quality: a more capable actor can
partially estimate $\Delta\vol(\Gcoop)$ from coalition structure signals,
shrinking the estimation error in Equation~\eqref{eq:sign_preservation} and
widening the margin of sign-correctness.

As these conditions strengthen, $\sign(\Deltahat(\pi))$ converges toward
$\sign(\Delta\vol(G) \mid \pi)$, and the local estimator approaches the true
objective. This gives GFM a self-improving property absent from frameworks with
fixed proxy objectives. In a fixed-utility framework, increased optimization
power amplifies Goodhart divergence, the agent becomes better at pursuing the
proxy without the proxy becoming more faithful to the underlying goal. The
local $\vol(G)$ estimator partially closes this gap because the actor is
optimizing and estimating the same quantity, and additional resources serve
both tasks simultaneously. A more capable GFM actor produces more accurate sign
estimates.

This anti-Goodhart property holds at the \emph{estimator} level---the fidelity
of $\Deltahat$ to $\Delta\vol(G)$---but does not by itself resolve the B-to-C
proxy failures identified in Section~\ref{sec:future_concerns}. A perfectly
accurate $\vol(G)$ estimator would still miss the Doll Problem and wireheading,
because those failures occur in the gap between capability volume and
experiential optionality, not in the gap between the estimator and the true
volume. Improved capability makes the actor better at maximizing $\vol(G)$;
whether the computational definition of $G$ captures enough of the experiential
structure (Candidate C) to avoid proxy divergence remains the open question
that the B/C framework raises. A capability space $\G$ that better approximates
experiential optionality would proportionally reduce the proxy's influence on
the objective.
